\documentclass{article}
\usepackage[
    a4paper,
    footskip = 1cm, 
    headsep = 0.3cm, 
    top = 2cm, 
    bottom = 2cm, 
    left = 2cm, 
    right = 2cm, 
]{geometry}

\usepackage{fontspec}  % using with LuaTeX to use cyrillic letters
\setmainfont{Times New Roman}
\setsansfont{Tahoma} % \textsf{your text}
\setmonofont{Consolas} % \texttt{your text} 
\usepackage[fontsize = 12pt]{fontsize} 

\usepackage{microtype} 
\usepackage[english, ukrainian]{babel}
\usepackage{amsmath}
\usepackage{nicefrac}
\usepackage{siunitx}
\usepackage{mathtools}
\usepackage[warnings-off={mathtools-colon, mathtools-overbracket}]{unicode-math}
\setmathfont{Latin Modern Math}
\usepackage{amsthm}
\usepackage{amsfonts}

\usepackage[
    unicode,
]{hyperref}
\hypersetup{
    colorlinks = true, 
    linkcolor = blue,
    breaklinks = false,	
    filecolor = magenta,
    citecolor = green,
    anchorcolor = black,
    urlcolor = cyan,
}

\newcounter{CurrentTask}
\setcounter{CurrentTask}{162}
\NewDocumentCommand{\Task}{}{
    \stepcounter{CurrentTask}
    \textbf{\theCurrentTask.}
}
\counterwithin{equation}{CurrentTask}

\begin{document}
\Task{}Розглядаючи передачу тепла не можна не зупинитися на явищі передачі за допомогою випромінювання. По суті процеси теплопередачі і енергообміну 
за допомогою випромінювання домінують у Всесвіті. Те що при сонячному затемненні або \textbf{??} разом зі світловим потоком зникає і тепловий, говорить 
про те, що тепло і світло поширюються з однаковою швидкістю, а саме зі швидкістю світла. Таким чином теплове випромінювання, має ту ж природу, що і світло, 
є по суті різновидом електромагнітного випромінювання подібно радіо або світлових хвиль. В силу своєї універсальності закони термодинаміки застосовані до 
будь-яких фізичних систем, в тому числі і до електромагнітного поля, що випромінюється тілом, що знаходиться в тепловій рівновазі з оточуючими тілами. Як 
показує повсякденний досвід при цьому необов'язково щоб тіла знаходилися в механічному контакті або були занурені в теплопровідну субстанцію типу рідини чи 
газу. Навіть вакуумноізольовані тіла, які обмінювалися тепловою енергією за допомогою випромінювання, з часом набувають однакову температуру. Для існування 
такої рівноваги, згідно нульового початку, усі три \textbf{???} тіло $1$, випромінювання і тіло $2$ повинні знаходиться в рівновазі один з одним. При такому 
випромінюванні, що знаходиться в рівновазі з випромінюючим тілом, говорять як про рівноважне випромінювання. Предметом нашого розгляду буде якраз таке 
випромінювання. Наші міркування, як завжди, будуть базуватися на експериментальних спостереженнях. Деякі положення, взяті з інших розділів фізики, 
здебільшого з оптики та електродинаміки, будуть застосовуватися без доведень. 

\Task{} Вперше ідею про променистий теплообмін в 1792 році висловив \textbf{???}. Ним же була висловлена думка про те, що кожне тіло випромінює енергію і в 
той же час поглинає ту ж саму енергію від інших тіл. При цьому більш нагріті тіла випускають більше енергії, ніж холодні тіла. Експерименти, проведені Леслі 
показали, що найбільш випромінювальною здатністю володіють зачернені поверхні, а найменшою - блискучі. У цих експериментах куб з гарячою водою, що має різні 
поверхні, що повертаються до реєструючого випромінювання приладу різними сторонами. Таким приладом може бути будь-який прилад, який реєструє теплове випромінювання, 
наприклад, болометр - відрізок платинової фольги, покритої комфорна \textbf{??? ???}, електроопір якої залежить від температури. Або термоелектрична батарея, 
що складається з безлічі послідовно з'єднаних термопар, гарячі спаї якої схильні до дії випромінювання, а холодні захищені. \vspace{5mm} \\
З іншого боку добре відомо, що зачернені тіла мають більшу поглинаючу здатність, ніж блискучі або просто світлі. Саме ці обставини спонукають жителів спекотних 
країн носити світлий одяг, а колби термостатів з цих же причин осріблюють. Звідси ми робимо висновок, що поверхні володіють високою випромінювальною здатністю, 
володіють і високою поглинальною здатністю і навпаки. Для додання цьому співвідношенню кількісного вмісту введемо деякі визначення Випромінювальною або 
променевипромінювальною здатністю називають величину рівну кількості енергії, яка випромінюється в одиницю часу з одиниці площі в одиницю тілесного кута - 
$I \left[\text{Вт}/(\text{м$^{2}$ ср})\right]$. Тоді енергія, яка переноситься випромінюванням в одиницю часу називається променистим потоком $d\Phi$ і дорівнює: 
\begin{equation}
    \label{eq:RadiantStream}
    d\Phi = I \cdot \cos{\theta} \cdot d\sigma \cdot d\Omega
\end{equation}
Де $\Theta$ кут між нормаллю до поверхні $d\Sigma$ і напрямом в якому випускається випромінювання. В оптиці для протяжних джерел використовують замість $I$ поняття 
поверхневої яскравості. Введемо також поняття світності як повного світлового потоку, який посилається одиницею поверхні в одну сторону тобто в тілесний кут $2\pi$. 
Якщо, як це має місце з тепловим випромінюванням, випромінювальна здатність не залежить від напрямку, тобто випромінювання ізотропне, то: 
\begin{equation}
\label{eq:IsotropicRadiation}
S = \int\limits_{2\pi}^{\textbf{??}} I \cdot \cos{\theta} \cdot d\Omega = 2\pi \int\limits_{0}^{\tfrac{\pi}{2}} I \cdot \cos{\theta} \cdot \sin{\theta} \, d\theta = \pi \cdot I
\end{equation}
Інша назва випромінювальної здатності - інтенсивність променистого потоку або інтенсивність випромінювання. Вона пов'язана з щільністю випромінювання $u$, тобто 
енергії випромінювання міститься в об'ємі простору. Через майданчик $d\sigma$ за час $dt$ проходить енергія $I = d\sigma dt d\Omega$. За цей час вона заповнить 
обсяг $c d\sigma dt$ заштрихований на малюнку. Тобто щільність променистої енергії складає:
\begin{equation}
    \label{eq:RadiantEnergyDensity}
    u = \frac{1}{c} \int I d\Omega = \frac{4\pi}{c} I
\end{equation}
Комбінуючи з \eqref{eq:IsotropicRadiation} знайдемо також:
\begin{equation}
    \label{eq:IsotropicRadiationCombined}
    S = \frac{c}{4}u
\end{equation}
Аналогічно вираженому для потоку частинок який проходить через одиничний  перетин в одиницю часу в газі. Аналогія не випадкова, якщо трактувати випромінювання, 
як газ фотонів які рухаються зі швидкістю $c$. \\
З точки зору термодинаміки випромінювання знаходиться в рівновазі з тілом, характеризується $T$, $p$ та $V$. З точки зору оптики випромінювання представляє 
суперпозицію електромагнітних хвиль всіляких частот. Наше завдання буде полягати в знаходженні зв'язку між макроскопічними параметрами, котрі характеризують 
випромінювання і його спектральними характеристиками, зокрема $I_{\omega}$ і $u_{\omega}$ записаними як:
\begin{equation}
    \label{eq:SpectralCharacteristics}
    u = \int u_{\omega} \, d\omega \text{ i } I = \int I_{\omega} \, d\omega
\end{equation}
Насамперед зазначимо, що в силу ізотропності нашого простору, щільність енергії випромінювання не залежить від обсягу. Таким чином, повна енергія 
випромінювання укладеного в деякому замкнутому об'ємі буде:
\begin{equation}
    \label{eq:TotalEnergyOfRadiation}
    U = u \cdot V
\end{equation}
Запишемо, що $u = u(t)$. Тут $T$ це температура випромінювання, за визначенням дорівнює температурі тіла, з яким випромінюванням знаходиться в рівновазі.

\Task{}Тепер покажемо, що щільність енергії випромінювання даної частоти $\omega$ покладену в деякий об’єм $V$ не залежить від матеріалу стінок або 
будь-яких тіл знаходяться в середині. Нехай випромінювання укладено між двома тілами $A$ і $B$ знаходяться при одній температурі $T$. І нехай для 
деякої частоти $\omega : u_{\omega}^{A} > u_{\omega}^{B}$. Встановимо між $A$ і $B$ фільтр пропускає тільки випромінювання частоти $\omega$. Тоді 
в силу \eqref{eq:IsotropicRadiationCombined} потік енергії з боку $A$ буде більше, ніж протилежний і тіло $B$ нагріється. Використовуючи 
отриману різницю температур для отримання роботи, тимчасово закривши абсолютно відображає екраном фільтр $A$. Дочекавшись поки між $A$ і $B$ 
встановиться однакова температура, повторимо процес, тим самим отримаємо вічний двигун другого роду. \\
Так як спектральна щільність uω не залежить від матеріалу стінок порожнини, в якій укладено випромінювання, а тільки від температури, то в 
силу \eqref{eq:RadiantEnergyDensity} інтенсивністю випромінювання частоти $\omega$ також є тільки формула температури. Тоді на площадку 
$d\sigma$ за час $dt$ в межах тілесного кута $d\Omega$, під кутом $\Theta$ до нормалі поверхні порожнини падає променистий потік \eqref{eq:RadiantStream}: 
\begin{equation}
    \label{eq:RadiantFlowToTheSquare_dsigma}
    d\Phi_{\omega} = I_{\omega} d\sigma \cdot dt \cdot d\Omega \cdot \cos{\theta} \cdot d\omega
\end{equation}
Частина його $(1 - \alpha_{\omega}) d\Phi_{\omega}$ відбивається. Тут $a_{\omega}$ - поглинаюча здатність стінок порожнини. На відбите 
випромінювання накладається власне випущених:
\begin{equation}
    \label{eq:ActuallyEmittedRadiation}
    E_{\omega} d\omega \cdot d\sigma \cdot dt \cdot \cos{\theta} \cdot d\Omega
\end{equation}
Таким чином в рівновазі: 
\begin{equation}
    \label{eq:BalanceEquation}
    (1 - \alpha_{\omega}) I_{\omega} + E_{\omega} = I_{\omega}
\end{equation}
Звідки отримуємо остаточно: 
\begin{equation}
    \label{eq:FinalCorrelation}
    \frac{E_{\omega}}{\alpha_{\omega}} = I_{\omega}
\end{equation}
Тут і в \eqref{eq:ActuallyEmittedRadiation} $E_{\omega}$ випромінююча здатність тіла порожнини на частоті $\omega$. Так як $I_{\omega}$ не залежить 
від матеріалу стінок порожнини, а тільки від $T$ і $\omega$ робимо висновок, що відношення випромінювальної здатності тіла до його поглинальної 
здатності є універсальна функція температури і частоти. Закон цей був висловлений Кірхгофом в 1859р. і носить його ім'я. Із закону випливає, що 
при даній температурі всяке тіло випромінює переважно такі промені, з такою частотою, з якою переважно поглинає. \\
Так, зелене скло, здається таким завдяки сильному поглинанню червоного світла, при нагріванні буде випускати червоне світло і навпаки. Добре 
приготовлене кварцове скло слабо поглинає світло будь-яких частот. Тому воно не світиться, будучи нагрітим аж до $1000\si{\celsius}$. З іншого боку, 
може здатися дивним, що оточуючі нас тіла виглядають різнокольоровими. Адже згідно \eqref{eq:FinalCorrelation} всі тіла повинні бути одного кольору. 
Це дійсно було б так, якби випромінювання, що йде від них, було б рівноважним. Однак це не так, оскільки реально такої рівноваги немає. Випромінення, 
що йде від Сонця близьке до рівноважного при температурі \textbf{???} Сонця. Поглинаючись в атмосфері, воно змінює свій спектральний склад і надалі 
переробляється освітленими тілами в інші нерівноважні спектри. Індивідуальність тіл позначається тут в повній мірі. У загальному випадку для будь-якого 
тіла (матеріалу) можна записати:
\begin{equation}
    \label{eq:RadiationEnergyBalanceRatio}
    \alpha_{\omega} + \rho_{\omega} + \tau_{\omega} = 1
\end{equation}
де $\rho_{\omega}$ і $\tau_{\omega}$ - відбивна і пропускна здатність речовини на частоті $\omega$. Очевидно, що \eqref{eq:RadiationEnergyBalanceRatio} є 
наслідком закону збереження енергії. \\
Тіло для якого $\alpha_{\omega} = 1$, а $\rho_{\omega} = \tau_{\omega} = 0$ для будь-яких частот $\omega$, називається абсолютно чорним тілом. Якщо $\rho_{\omega} = 1$, 
то тіло називають абсолютно білим. Якщо, $\tau_{\omega} = 1$, а $\alpha_{\omega} = \rho_{\omega} = 0$ то абсолютно прозорим. Фундаментальне значення має поняття абсолютно 
чорного тіла, тому що згідно з \eqref{eq:FinalCorrelation} спектр випромінювання такого тіла і є спектр рівноважного випромінювання при температурі $T$. Зрозуміло таке 
тіло є математичною абстракцією. Для реальних тіл: $\alpha_{\omega} \sim 1$; так сажа має $\alpha_{\omega} = 1$, ще вище $\alpha_{\omega}$ у платинової черні. Майже ідеальною 
є спектр випромінювання виходить з невеликого отвору зробленого в замкнутій порожнині нагрітої до температури $T$. Випромінювання потрапляє всередину такої порожнини практично 
націло поглинається її стінками в результаті багаторазових відображень. Саме такого роду конструкції використовуються для проведення особливо точних експериментів. Важливу роль 
в теорії грають абсолютно білі тіла. На перший погляд може здатися, що випромінювання температури $T_{1}$, що потрапило в порожнину з абсолютною відбивають стінками при температурі 
$T_{1}$ зможе перебувати там необмежено \textbf{???} ставлячи під сумнів справедливість нульового початку. Однак така ситуація не буде стійкою. Як вперше помітив Планк навіть, 
незначна пилинка рано чи пізно перетворить нерівноважне випромінювання в рівноважне. Доведення цього розглянемо пізніше.

\Task{} Звернемося тепер до питання про тиск випромінювання. Вперше на можливість тиску, створеного випромінюванням, вказав Кеплер, помітивши, що хвіст комет завжди спрямований в бік 
від Сонця. Насправді це явище в великій мірі обумовленого тиском сонячного вітру, на поставлене питання знайшов своє рішення в 1897р. в \textbf{???} експерименті А.\textbf{???}. Нехай 
випромінювання заповнює циліндр з абсолютно білими стінами, торці якого заповнені що поглинає матеріалом при температурі $T_{1}$ і $T_{2}$. Заберемо абсолютно білу заслонку $p_{2}$. 
Порожнина заповниться випромінюванням при $T_{2}$. Вставимо тепер $p_{2}$ заберемо $p_{1}$ і передвинемо $p_{2}$ до $A$. Тіло $A$ поглине випромінювання з температурою $T_{2}$ і нагріється. 
Вставимо тепер $p_{1}$ на місце $p_{2}$ і повторимо процес. Таким чином можна передати тепло від $B$ до $A$ навіть якщо $T_{1} > T_{2}$. Але на другому початку це можливо тільки якщо 
виконується робота по переміщенню заслінок проти сил тиску, створюваного випромінюванням. Тобто випромінювання створює тиск. Раніше, на підставі молекулярно-кінетичних уявлень було показано, що 
\begin{equation}
    \label{eq:PressureBasedOnMolecular-KineticConcepts}
    p = \frac{1}{3} u
\end{equation}
При виведенні була використана ізотропність випромінювання, що вірно для теплового випромінювання, а також що відображення від стінок є абсолютно пружним, тобто має місце для білого тіла. 
Навіть якщо це не так, то в рівновазі на кожен промінь, що падає за час $dt$ на майданчик $d\sigma$ під кутом $\Theta$ завжди знайдеться промінь випущений під кутом $\Theta$ з цього ж майданчика. 
Зауважимо принагідно, що газ фотонів швидше подібний не до ідеально газу, а до насиченої пари. Справді той факт, що щільність енергії випромінювання залежить від температури, говорить про те, 
що квантовий газ не пружинить подібно насиченій парі. Точно така ж формула \eqref{eq:PressureBasedOnMolecular-KineticConcepts} для ізотропного випромінювання була отримана методами електродинаміки 
Максвела, а експериментальне підтвердження існування тиску світла було отримано Лебедєвим П. І. в 1901р., а пізніше Столєтовим.

\Task{}Скористаємося тепер основним наслідком другого початку і застосуємо його до рівноважного виміру. Тоді:
\begin{equation}
    \label{eq:MainConsequenceOnEquilibriumDimension_step1}
    \left(\frac{\partial(uV)}{\partial V}\right)_{T} = u = T \left(\frac{\partial p}{\partial T}\right)_{V} - p = 
    \frac{T}{3} \left(\frac{\partial u}{\partial T}\right)_{V} - \frac{1}{3} u
\end{equation}
Звідси
\begin{equation}
    \label{eq:MainConsequenceOnEquilibriumDimension_step2}
    \frac{du}{dT} = 4 \frac{u}{T}
\end{equation}
Або
\begin{equation}
    \label{eq:EnergyValueFromMainConsequenceOnEquilibriumDimension}
    u = \alpha T^{4}
\end{equation}
Саме таким чином Больцман в 1884р. встановив теоретично залежність \eqref{eq:EnergyValueFromMainConsequenceOnEquilibriumDimension}. Раніше Стефан в 1879р. встановив цю залежність експериментально, 
де $\alpha = 7.64 \cdot 10^{-16}$ Дж/(м$^{3} \cdot $ K$^{4}$). Формула \eqref{eq:EnergyValueFromMainConsequenceOnEquilibriumDimension} може бути отримана з тих міркувань, що ентропія - є функція стану, 
а також методом циклів. Знайдемо ентропію випромінювання: 
\begin{equation}
    \label{eq:RadiationEntropy}
    dS = \frac{d(uV)}{T} + \frac{1}{3} \frac{u}{T} dV = \frac{V}{T} dU + \frac{4}{3} dV = 4 \alpha V T^{2} dT + \frac{4}{3} \alpha T^{3} dV
\end{equation}
Або:
\begin{equation}
    \label{eq:RadiationEntropy_versionTWO}
    S = \frac{4}{3} \alpha T^{3} dV
\end{equation}
Звідси отримаємо рівняння адіабати:
\begin{equation}
    \label{eq:AdiabaticEquation}
    T^{3} V = const \text{ або } pV^{\tfrac{4}{3}} = const
\end{equation}
Відповідно постійна адіабати $\gamma = 4/3$. \textbf{???}, що наближаючись \textbf{???}, можна зробити висновок, що кількість \textbf{???}. Однак цей висновок не є правомірним. По-перше, для газу фотонів 
закон рівнорозподілу енергії за ступенями свободи повинен застосовуватися з деякими застереженнями, так як фотони є істотно квантові об'єкти. По-друге, так як для випромінювання $p = p(t)$, то \textbf{???} 
$C_{p} = C_{t} - \infty$. При цьому:
\begin{equation}
    \label{eq:HeatCapacityAtConstantVolume}
    C_{V} = T \left(\frac{\partial S}{\partial T}\right)_{V} = 4 \alpha T^{2} V
\end{equation}
Внутрішня енергія випромінювання $U$ дорівнює:
\begin{equation}
    \label{eq:InternalEnergyOfRadiation}
    U = uV = \alpha T^{4} V
\end{equation}
Ентальпія:
\begin{equation}
    \label{eq:Enthalpy}
    H = U + pV = \alpha T^{4} V + \frac{1}{3} \alpha T^{4} V = \frac{4}{3} \alpha T^{4} V
\end{equation}
Вільна енергія Гельмгольца:
\begin{equation}
    \label{eq:HelmholtzFreeEnergy}
    F = U - TS = \alpha T^{4} V - \frac{4}{3} \alpha T^{4} V = -\frac{1}{3} \alpha T^{4} V
\end{equation}
Нарешті вільна енергія Гібса:
\begin{equation}
    \label{eq:GibbsFreeEnergy}
    G = F + pV = -\frac{1}{3} \alpha T^{4} V + \frac{1}{3} \alpha T^{4} V = 0
\end{equation}
Той факт, що $G = 0$ говорить про те, що $G$ не може бути характеристичною функцією для випромінювання, так як $p = p(T)$, подібно до того, як це має місце з \textbf{???} \textbf{???} . Крім того з \eqref{eq:GibbsFreeEnergy} 
випливає, що $\mu = \left(\frac{\partial G}{\partial N}\right)_{p, T} = 0$ тобто \textbf{???} фотонного газу дорівнює $0$. Іншими словами утворення або зникнення фотонів не сполучене з виконанням роботи. 

\Task{}Термодинаміка, даючи отримані важливі співвідношення особисто не може дати відносно точного виду спектральної щільності. Зазвичай, в термодинаміці вдається отримати деякі важливі відомості щодо виду функції $u_{\omega}$ та 
$I_{\omega}$. \vspace{5mm} \\
Розглянемо сферичну порожнину з абсолютно білими стінками, які рівномірно розширюються зі швидкістю $r$. Чи буде спочатку рівноважне випромінювання з температурою $T$, залишатися рівноважним і далі? Адже спектральний склад 
випромінювання в силу \textbf{???} буде змінюватися і неочевидно, що його вид буде відповідати деякому рівноважному випромінюванню з $T'$. Введемо в оболонку нескінченно малу абсолютно поглинаючу пилинку. З плином часу пилинка 
перетворить навіть нерівноважне випромінювання в рівноважне. Цей процес йде мимовільно і є незворотнім. Коли випромінювання стане рівноважним, не забираючи пилинку почнемо стискати оболонку і доведемо її об’єм до початкового. 
Енергія пилинки мала, і мало впливає на загальну енергію випромінювання. Оскільки тиск ізотропного випромінювання залежить від повної щільності \eqref{eq:PressureBasedOnMolecular-KineticConcepts} випромінювання, а не від його 
спектрального складу, то робота випромінювання при розширенні оболонки буде рівна роботі при стисканні. Звідси випливає, що енергія випромінювання в результаті адіабатичного розширення з наступним стисканням не змінюється, а 
отже не змінюється температура і тиск в такому круговому процесі. Оскільки робота не виконалася тепло не передавалося, то в оточуючих тілах не виникло ніяких змін і, як результат процес є оборотним. Проте за припущенням одна 
із стадій не оборотна. Це доводить, що припущення про нерівноважність випромінювання при адіабатичному розширенні не вірно. В 1894 Вінн довів, що рівноважне випромінювання замкнене в оболонці з ідеально відбиваючими стінками, 
буде залишатися рівноважним при квазістатичному стисканні чи розширенні. Тепер звернемося до рисунка. Час між двома послідовними відбиттями складає \textbf{???}. За цей час радіус оболонки змінюється. Зміна частоти через ефект 
Доплера складає:
\begin{equation}
    \label{eq:FrequencyShift_due_to_DopplerEffect}
    \Delta \omega \tfrac{\Delta \omega}{\omega} = -2\frac{r \cos{\theta}}{c} = - 2 \frac{\Delta r \cos{\theta}}{c \Delta t} = - \frac{\Delta r}{r}
\end{equation}
Інтегруючи отримаємо:
\begin{equation}
    \label{eq:FrequencyShift_Integrated}
    \omega r = const
\end{equation}
Врахувавши те, що $V \sim r^{3}$:
\begin{equation}
    \label{eq:Consider_to_V_Approximation}
    V \omega^{3} = const
\end{equation}
Скориставшись адіабатичністю процесу отримаємо:
\begin{equation}
    \label{eq:WienDisplacementLaw}
    \frac{\omega}{T} = const
\end{equation}
Іншими словами величина $\frac{\omega}{T}$ являється адіабатичним варіантом рівноважного випромінювання. На перший погляд \eqref{eq:WienDisplacementLaw} має частотний характер, а не характеризує випромінювання. Припустимо, що 
це не так і тоді в якому-небудь процес отримаємо, що $\omega = f(T)$ і \eqref{eq:WienDisplacementLaw} не виконується. Змінимо адіабатичну температуру випромінювання від $T_{1}$ до $T_{2}$, далі з допомогою названого процесу 
з $T_{1}$ до $T_{2}$. При цьому спектральний склад випромінювання змінюється, а температура ні, що заперечує розвиненим уявленням про універсальний характер спектральної щільності рівноважного випромінювання. 

\Task{}Закон Стефана-Больцмана про пропорційність щільності енергії випромінювання $T^{4}$ залишає без відповіді питання про вид функції $W\omega$ - спектральної щільності енергії випромінювання. Скористаємось тим, що для 
будь-якого квазістатичного процесу $\omega/T$ є інваріанта. Введемо нову змінну. Тоді приймаючи $\xi = \omega/T$ отримаємо: 
\begin{equation}
    \label{eq:TotalEnergyDensity_of_blackbodyConsequence_by_Stefan-BoltzmannLaw}
    u = \int\limits_{0}^{\infty} u_{\omega} \, d\omega = \int\limits_{0}^{\infty} u_{\omega}(\omega, T) \, d\omega = 
    \int\limits_{0}^{\infty} u(\xi) T \, d\xi = \alpha T^{4}
\end{equation}
Виконання \eqref{eq:TotalEnergyDensity_of_blackbodyConsequence_by_Stefan-BoltzmannLaw} можливе, якщо $u_{\omega}(\omega, T)$ має спеціальний вид: 
\begin{equation}
    \label{eq:SpecialType_of_u_Function}
    u_{\omega}(\omega, T) = \omega^{3} f(\frac{\omega}{T})
\end{equation}
Тоді
\begin{equation}
    \label{eq:TotalEnergyDensity_of_blackbodyConsequence_by_Stefan-BoltzmannLawFinal}
    u = \int\limits_{0}^{\infty} u_{\omega} \, d\omega = T^{4} \int\limits_{0}^{\infty} \xi^{3} f(\xi) \, d\xi
\end{equation}
де, по-перше, $f(\xi)$ має бути таким, щоб інтеграл сходився, а, по-друге величина $f(\xi)$ - розмірна з розмірністю $E \cdot \tau^{4} \cdot l^{-3}$. Тут: $E$ - енергія, $\tau$ - час, $l$ - довжина. Коли мова йде про випромінювання 
в $f(\omega/T)$ можуть входити частота, температура і швидкість світла. Оскільки фотони безмасові частинки, то у вираз для $f$ не можуть входити ніякі величина з розмірністю маси (типу маси електрона). З точки зору класичної фізики 
єдиною можливістю організувати величину з розмірністю енергії залишається використання сталої Больцмана в комбінації  з температурою. \vspace{5mm} \\
Тоді отримаємо, що $f(\omega, T) = \frac{k_{\text{Б}} \cdot T}{c^{3} \cdot \omega} \varphi(\tfrac{\omega}{T})$ чи: 
\begin{equation}
    \label{eq:SpectralDensity_of_Radiation_Energy}
    u_{\omega}(\omega, T) = \frac{k_{\text{Б}} \cdot T \cdot \omega^{2}}{c^{3}} \varphi(\tfrac{\omega}{T})
\end{equation}
де $\varphi(\tfrac{\omega}{T})$ вже функція безрозмірна. Формула \eqref{eq:SpecialType_of_u_Function} отримала назву закону зміщення Віна для структурної функції, встановлена Віном в 1893 році. В деякому іншому вигляді, з врахуванням 
того, що $\lambda = \frac{c}{\omega}$, цей закон записується як:
\begin{equation}
    \label{eq:SpectralDensity_of_Radiation_Energy_Other_Appearance}
    u_{\lambda}(\lambda, T) = \lambda^{-5} f(\lambda, T)
\end{equation}
Зауважимо, що оскільки інтеграл в \eqref{eq:TotalEnergyDensity_of_blackbodyConsequence_by_Stefan-BoltzmannLawFinal} має сходитися, підінтегральний вираз, а, слідуючи і uω є функція з максимумом. Дискретизуючи 
\eqref{eq:SpectralDensity_of_Radiation_Energy_Other_Appearance} отримаємо, що:
\begin{equation}
    \label{eq:SpectralDensity_of_Radiation_Energy_Other_Appearance_Discretized}
    -5 \lambda_{m}^{-6} f(\lambda_{m}, T) + \lambda_{m}^{-5} \cdot T \cdot f(\lambda_{m}, T) = 0 
\end{equation}
де $\lambda_{m}$ - довжина хвилі, яка відповідає максимуму спектральної щільності рівноважного випромінювання. \vspace{5mm} \\
Звідси отримаємо, що:
\begin{equation}
    \label{eq:WienDisplacementLaw_versionTWO}
    \lambda_{m} T = \frac{5 f(\lambda_{m}, T)}{f'(\lambda_{m}, T)} = b = const
\end{equation}
Вимірювання дають значення $b = 0.2898$ см $\cdot$ К. Формула \eqref{eq:WienDisplacementLaw_versionTWO} дає закон зміщення Віна, згідно якого довжина хвиль, яка відповідає максимуму спектральної щільності 
рівноважного випромінювання обернено пропорційна абсолютній температурі. Відповідно \eqref{eq:WienDisplacementLaw_versionTWO} може бути переписана у вигляді: 
\begin{equation}
    \label{eq:WienDisplacementLaw_versionTWO_Other_Appearance}
    \frac{\omega_{m}}{T} = const
\end{equation}
де взагалі говорячи максимуми для $\omega$ і $\lambda$ в спектральній щільності не відповідають один одному. Відношення \eqref{eq:WienDisplacementLaw_versionTWO} і \eqref{eq:WienDisplacementLaw_versionTWO_Other_Appearance} 
використовується в пірометрах - приладах для вимірювання високих температур, відповідно до максимуму спектральної щільності, використаних в діапазоні вище $1000\si{\celsius}$ . При цьому використовують термін - кольорова температура.

\Task{}
Отримані раніше термодинамічні відношення, закони Стефана-Больцмана і Віна, структурна формула закону Віна для структурної функції $f(\omega/T)$ це напевно все, що може дати термодинаміка базуючись на класичних уявленнях про 
природу випромінювання. Крім того детальний аналіз виразу для $u_{\omega}$ показує на його внутрішнє протиріччя. Насамперед безрозмірна функція $\varphi(\omega/T)$, яка входить до \eqref{eq:SpectralDensity_of_Radiation_Energy} 
містить в якості аргументу розмірне відношення. Класична фізика не вказує яким чином можна \textbf{???} виміряти це відношення так, щоб в безрозмірну структурну функцію входив безрозмірний аргумент. Далі помітимо, що при малих 
частотах $\varphi(\omega/T)$ може бути розкладена в ряд по малому параметру $\omega/T$: 
\begin{equation}
    \label{eq:DimensionlessFunctionExpandedIntoSeries}
    \varphi\left(\frac{\omega}{T}\right) = \varphi(0) + \varphi'(0) \tfrac{\omega}{T} + \ldots
\end{equation}
Причому $\varphi(0) \neq 0$. В іншому випадку $u_{\omega}$ не залежало від температури в границі малих $\omega\left(u_{\omega} = \frac{k_{\text{ Б }} \cdot \omega^{3}}{c^{3}} \varphi'(0)\right)$. Тоді два тіла, які знаходяться 
при різних температурах і розділені фільтром, яких пропускає тільки низькочастотне випромінювання ніколи не досягли б рівноважності. Звідси знайдемо, що при $\omega \to 0$: 
\begin{equation}
    \label{Approximation_of_SpectralDensity_at_omega_to_0}
    u_{\omega} = A \frac{k_{\text{ Б }} \cdot T \cdot \omega^{2}}{c^{2}}, \text{ де } A = \varphi(0)
\end{equation}
Саме до такої залежності $u_{\omega}(\omega, T)$ прийшли в 1900 році Реле і Джинс, помітивши теорему про рівнорозділення за степенями волі до рівноважного випромінювання у всьому частотному діапазоні. Але в такому випадку не 
задовольняє закон Стефана-Больцмана. Для розуміння суті протиріччя повторимо роздуми Релея і Джинса. Згідно теореми на кожну коливальну степінь волі приходиться енергія $k_{\text{ Б }}T$. Отже, все, що потрібно - це знайти 
кількість степенів свободи випромінювання в порожнині об’єму $V$. Оскільки рівноважне випромінювання не залежить від форми порожнини, можна вибрати найбільш зручну для розрахунків, а саме куб зі стороною $L$. Будь-яке можливе 
розподілення електромагнітного поля можна представити у вигляді суперпозиції стоячих хвиль (подібно коливанням струни) утворюючих повну систему ортогональних функцій. Розрахуємо кількість стоячих хвиль, які вкладаються в кубі з 
частотами від $\omega$ до $\omega + d\omega$. Це і буде кількістю степеней свободи, яке необхідне для описання будь-якої структури електромагнітного поля з частотою $\omega$, $\omega + d\omega$. У кожному з напрямків довжини хвиль, 
які вкладаються в одному напрямку утворюють послідовність $2L$, $L$, $3L$, $\dots$. Або відповідають послідовності хвильових векторів стоячих хвиль $k_{x}, k_{y}, k_{z} = \tfrac{\pi}{L}, \tfrac{2\pi}{L}, \tfrac{3\pi}{L}, \dots$. 
У фазовому просторі $k_{x}, k_{y}, k_{z}$, кожному набору  $\tilde{k}_{x}, \tilde{k}_{y}, \tilde{k}_{z}$ відповідає точка. Точки ці утворюють правильну решітку з кубиків з довжиною ребра $\pi/L$ і  об’ємом $(\pi/L)^{3}$ \textbf{???}. 
Об’єм частини шару радіусу рівний одному додатному октанту дорівнює $\tfrac{1}{8} \left(\frac{4\pi}{3}\right)k^{3} = \left(\frac{\pi}{6}\right)k^{3}$. Для досить великих $\omega$, величина $k = \omega / c$ велика по відношенню до 
ребра елементарного кубика $\pi/L$. Тому число стоячих хвиль, які попадають в октант, дорівнює числу степеней свободи, які проходять на випромінювання з частотою від $0$ до $\omega$ складає: 
\begin{equation}
    \label{eq:AmountDegrees_of_Freedom}
    Z = \frac{\tfrac{\pi}{6} k^{3}}{\left(\tfrac{\pi}{L}\right)^{3}} = \frac{V}{6 \pi^{2}} \cdot k^{3} = \frac{V}{6 \pi^{2}} \cdot \frac{\omega^{3}}{c^{3}}
\end{equation}
Врахуємо, що на кожну із хвиль припадає дві поляризації. Тоді на інтервал $d\omega$ кількість степеней волі складає:
\begin{equation}
    \label{eq:AmountDegrees_of_Freedom_per_Interval}
    dz = z \tfrac{dz}{d\omega} d\omega = \tfrac{V}{\pi^{2}} \cdot \tfrac{\omega^{2}}{c^{3}} d\omega
\end{equation}                                              
Тоді на цей інтервал $d\omega$ приходиться енергія $d\varepsilon = dz \cdot kT$ зі щільністю:
\begin{equation}
    \label{eq:EnergyDensity_on_Interval}
    u_{\omega} = \frac{k_{\text{Б}} \cdot T}{\pi^{2}} \cdot \frac{\omega^{2}}{c^{3}}
\end{equation}
Отримане співвідношення називається формулою Релея-Джинса. Вона досить точно описує експериментальні дані в радіо і далекому інфрачервоному діапазоні. Проте її використання по всьому діапазоні частот дає не суттєвий результат - 
інтегральна щільність енергії дорівнює нескінченності:
\begin{equation}
    \label{eq:IntegralEnergyDensity}
    u = \int\limits_{0}^{\infty} u_{\omega} \, d\omega = \frac{k_{\text{Б}} \cdot T}{\pi^{2}} \cdot \frac{1}{c^{3}} \cdot \int\limits_{0}^{\infty} \omega^{2} \, d\omega = \infty
\end{equation}
У свій час цей результат так вразив сучасників, що по місцевому виразу \textbf{???}, його назвали ультрафіолетовою катастрофою. ЇЇ причиною можна назвати те, що речовина має скінченну кількість степенів свободи, 
а випромінювання нескінчене. Таким чином теорема про рівнорозподіл енергії по степеням волі не працює у класичному виді у випадку рівномірного випромінювання і речовини.

\Task{}Виклик кинутий природою прийняли провідні фізики \textbf{?????}. Проте кінцевий розв’язок було знайдено Планком, який спочатку емпірично, а потім теоретично 14 грудня 1900 року на засіданні Берлінського фізичного 
товариства обґрунтував формулу для спектральної щільності випромінювання. Планк висловив гіпотезу, згідно за якою випромінювання випускається квантами енергії, які можуть приймати дискретні значення $\varepsilon_{0}$. Ці кванти 
випромінюються і поглинаються молекулами речовини, які можуть бути представлені як гармонічні осцилятори, енергія яких може приймати лише дискретні значення $0, \varepsilon, 2\varepsilon, 3\varepsilon, \dots$. В стані 
термодинамічної рівноваги між актами поглинання і випромінювання установлюється детальна рівновага, в результаті якої всі рівні гармонічного осцилятора збуджуються, але з різною ймовірністю. Тут під осцилятором потрібно 
розуміти не тільки частинки, а й стоячі хвилі в порожнині. Ймовірності, з якими збуджуються хвилі, пропорційні згідно формули Максвела-Больцмана: 
\begin{equation}
    \label{eq:Maxwell-BoltzmannFormula}
    e^{\frac{\varepsilon}{k T}}
\end{equation}
Тоді середня енергія осцилятора дорівнює:
\begin{equation}
    \label{eq:AverageOscillatorEnergy}
    \langle \mathcal{E} \rangle = \frac{\sum\limits_{n=0}^{\infty}!n \varepsilon \cdot \exp\left\{-\tfrac{n \cdot \varepsilon}{k_{\text{Б}} \cdot T}\right\}}{\sum\limits_{n=0}^{\infty}! \varepsilon \cdot \exp\left\{-\tfrac{n \cdot \varepsilon}{k_{\text{Б}} \cdot T}\right\}} =
    \frac{\sum\limits_{n=0}^{\infty}!n \varepsilon \cdot \exp\left\{-\beta \cdot n \cdot \varepsilon\right\}}{\sum\limits_{n=0}^{\infty}!\varepsilon \cdot \exp\left\{-\beta \cdot n \cdot \varepsilon\right\}}
\end{equation}
де як завжди $\beta = \tfrac{1}{kT}$. Помітимо, що чисельники отримуються диференціюванням знаменника по $-\beta$. \vspace{5mm} \\
Знаменник, будучи геометричною прогресією дорівнює:
\begin{equation}
    \label{eq:Denominator_of_AverageOscillatorEnergy}
    \sum\limits_{n=0}^{\infty} ! e^{-\beta \cdot n \cdot \varepsilon} = \frac{1}{1 - e^{-\beta \cdot \varepsilon}}
\end{equation}
Диференціюючи \eqref{eq:Denominator_of_AverageOscillatorEnergy}, знаходимо:
\begin{equation}
    \label{eq:AverageOscillatorEnergyValue}
    \langle \mathcal{E} \rangle = \frac{\varepsilon}{e^\tfrac{\varepsilon}{kT} - 1}
\end{equation}
Підставляючи в \eqref{eq:AverageOscillatorEnergyValue} в \eqref{eq:AmountDegrees_of_Freedom_per_Interval} нарешті знайдемо:
\begin{equation}
    \label{eq:PlanckFormula}
    u_{\omega}(\omega, T) = \frac{\omega^{2}}{\pi^{2} \cdot c^{3}} \cdot \frac{\varepsilon}{e^\tfrac{\varepsilon}{kT} - 1} = \frac{\hslash \omega^{3}}{\pi^{2} c^{3}} \cdot 
    \frac{1}{e^\tfrac{\hslash \omega}{kT} - 1}
\end{equation}                                       
Отримана формула називається формулою Планка. Помітимо, що $\varepsilon$ може бути функцією тільки частоти самого осцилятора і не залежить від $T$. В такому випадку для того, щоб аргумент функції $\varphi(\omega/T)$ був безрозмірним 
необхідно і достатньо, щоб:
\begin{equation}
    \label{eq:Value_of_epsilon}
    \varepsilon =\hslash \cdot \omega
\end{equation}
де $\hslash$ - універсальна стала Планка, яка дорівнює $\hslash = 1.0546 \cdot 10^{-27}$ ерг*с. Вона універсальна, оскільки універсальна функція $u_{\omega}(\omega, T)$ є функцією  безрозмірного відношення $\hslash \omega / kT$. \vspace{5mm} \\
При малих частотах при виконанні умови \eqref{eq:PlanckFormula} переходить в $\hslash \omega << kT$: 
\begin{equation*}
    \label{eq:Rayleigh-JeansFormula}
    u_{\omega}(\omega, T) = \frac{\hslash \omega^{3}}{\pi^{2} \cdot c^{3}} \cdot \frac{1}{1 + \frac{\hslash \omega}{kT} - 1} = \frac{\omega^{2}}{\pi^{2} c^{3}}
\end{equation*}
Тепер ми можемо уточнити значення сталих, які входять в формули Стефана-Больцмана і Віна. Інтегруючи \eqref{eq:PlanckFormula} отримаємо: 
\begin{equation}
    \label{eq:FINALValue_of_IntegralEnergyDensity}
    u = \frac{\hslash}{\pi^{2} \cdot c^{3}} \int\limits_{0}^{\infty} \frac{\omega^{3} d\omega}{e^{\tfrac{\hslash \omega}{kT}} - 1} = \frac{k_{\text{Б}}^{4} \cdot T^{4}}{\pi^{2} \cdot c^{3} \cdot \hslash^{3}} 
    \int\limits_{0}^{\infty} \frac{\xi^{3} d\xi}{e^{\xi} - 1} = \frac{\pi^{2}}{15} \cdot \frac{(k_{\text{Б}}\cdot T)^{4}}{c^{3} \hslash^{3}}
\end{equation}
На практиці використовується вираз для світимості $S$, яка \eqref{eq:IsotropicRadiation}, \eqref{eq:RadiantEnergyDensity} дорівнює: 
\begin{equation}
    \label{eq:LightingFormula}
    S = \frac{c}{4} u = \frac{\pi^{2}}{60} \cdot \frac{k_{\text{Б}}^{4}}{c^{3} \hslash^{3}} \cdot T^{4} = \sigma \cdot T^{4}
\end{equation}
де $\sigma = 5.67 \cdot 10^{-8}$ Вт$\cdot$м$^{-2}\cdot$К$^{-4}$ стала Стефана-Больцмана.
Помітимо також, що випромінююча здатність реальних тіл визначається як відношення інтенсивності випромінювання (або світимості) тіла до інтенсивності світимості абсолютно чорного тіла. Якщо це відношення близьке до одиниці, то тіло називають сірим. 
При цьому зазвичай використовується закон, аналогічний Стефану-Больцману в формі $S = \sigma'\cdot T^{n}$, де $n$ і $\sigma'$ визначаються емпірично (Люмер). Найбільш точний вираз для інтерпретації експериментальних даних запропонований Гельфготом, 
згідно якому: 
\begin{equation}
    \label{eq:HelfgottFormula_for_Lighting}
    S = \left(1 - e^{-\alpha T}\right) \sigma T^{4}
\end{equation}
де $\alpha$ - емпірична стала. З \eqref{eq:HelfgottFormula_for_Lighting} видно, що чим вища температура, тим ближче властивості випромінювання до властивостей рівноважного випромінювання.
\end{document}